%% kapa: exam=1 round=34 year=2023 period=1 form=A subject=civil_law subject_ko=민법 seq=16 number=16
%% source: Q-Net 감정평가사 기출문제 (cst00309) artlSeq=5211789 file=raw/1008/5211789/2023년도 제34회 감정평가사 1차시험 1교시.pdf
%% conversion: kordoc 4.24.51 -> tools/gen1.py (auto); review=proofread-sonnet; flags=none
\begin{question}{16}
법률행위의 무효에 관한 설명으로 옳지 않은 것은? (다툼이 있으면 판례에 따름)
\begin{choices}
  \item 무권대리행위에 대한 본인의 추인은 다른 의사표시가 없는 한 소급효를 가진다.
  \item 법률행위의 일부분이 무효일 때, 그 나머지 부분의 유효성을 판단함에 있어 나머지 부분을 유효로 하려는 당사자의 가정적 의사를 고려하여야 한다.
  \item 토지거래허가구역 내의 토지를 매매한 당사자가 계약체결시부터 허가를 잠탈할 의도였더라도, 그 후 해당 토지에 대한 허가구역 지정이 해제되었다면 위 매매계약은 유효가 된다.
  \item 무효인 법률행위를 추인에 의하여 새로운 법률행위로 보기 위해서는 당사자가 그 무효를 알고서 추인하여야 한다.
  \item 처분권자는 명문의 규정이 없더라도 처분권 없는 자의 처분행위를 추인하여 이를 유효하게 할 수 있다.
\end{choices}
\end{question}
\answer{3}
